\section{Datadictionary landbouwgebruikspercelen}
\label{app:datadictionary}

\setlength{\LTpre}{0pt}
\setlength{\LTpost}{0pt}

{
    \footnotesize
    \begin{xltabular}{\textwidth}{@{} >{\raggedright\arraybackslash}l >{\raggedright\arraybackslash}X @{}}
        \caption{Beschrijving van de attributen in de landbouwgebruikspercelen dataset \autocite{vlaanderen_landbouwpercelen_2025}}\label{app:ontwikkelingTool_landbouwkaart-vlaanderen-attributen} \\

        \toprule
        \textbf{Attribuut} & \textbf{Omschrijving} \\
        \midrule
        \endfirsthead


        AGPAKEY & Unieke identificatiesleutel van het landbouwgebruiksperceel binnen één campagne \\
        BT\_OMSCH & Bedrijfstypologie (economische specialisatie) \\
        BT\_BRON & Bron van de bedrijfstypologie (jaartal berekening of aangegeven specialisatie) \\
        PRC\_NMR & Perceelsnummer van het perceel in de verzamelaanvraag \\
        GRAF\_OPP & Grafische oppervlakte (ha, tot 1m² nauwkeurig) \\
        REF\_ID & Europees verplicht perceelsID. Verandert mee met de verschijningsvorm van het perceel \\
        GWSCOD\_V & Code voorteelt \\
        GWSNAM\_V & Naam voorteelt \\
        GWSCOD\_H & Code hoofdteelt \\
        GWSNAM\_H & Naam hoofdteelt \\
        GWSGRPH\_LB & Naam groep hoofdteelt \\
        GWSCOD\_N & Code eerste nateelt \\
        CWSNAM\_N & Naam eerste nateelt \\
        GWSCOD\_N2 & Code tweede nateelt \\
        GWSNAM\_N2 & Naam tweede nateelt \\
        AMKM & Code agromilieumaatregel \\
        AMKM\_LB & Naam agromilieumaatregel \\
        ECOREGELING (ECOR) & Code ecoregeling \\
        ECOREGELING\_LB & Naam ecoregeling \\
        BLS & Code aanplantsubsidie boslandbouwsystemen \\
        BLS\_LB & Naam aanplantsubsidie boslandbouwsystemen \\
        GESP\_PM & Gespecialiseerde productiemethode \\
        GESP\_PM\_LB & Beschrijving gespecialiseerde productiemethode \\
        BOPAKCOD & Code beheerovereenkomst VLM \\
        BOPAKCOD\_L(B) & Naam beheerovereenkomst VLM \\
        BIOCERT & Perceel onder bio-controle bij een bio-controleorgaan (J/N) \\
        ERO\_NAM & Erosiecode (kleur) van het perceel \\
        STAT\_BGV & Status Blijvend Grasland (BG) \\
        MEERJARIG\_GRASLAND (MG) & Status Meerjarig Grasland (MG6 of hoger). Voorbeeld: “MG16 = 16de jaar grasland” \\
        LANDBSTR & Landbouwstreek waarin het middelpunt van het perceel ligt \\
        STAT\_AAR & Status Aardappelen, opvolgen rotatieplicht \\
        PCT\_EKBG & Percentage van het perceel dat ecologisch kwetsbaar blijvend grasland is (range) \\
        PCT\_WETVEE(N) & Percentage van het perceel dat wetland en/of veengebied is (range) \\
        PRC\_GEM & Gemeente waarin het middelpunt van het perceel zich bevindt \\
        PRC\_NIS & NIS-code van de gemeente waarin het middelpunt van het perceel zich bevindt \\
        X\_REF & X-coördinaat van het middelpunt van het perceel (Lambert) \\
        Y\_REF & Y-coördinaat van het middelpunt van het perceel (Lambert) \\
        WGS84\_LG & Lengtegraad van het middelpunt van het perceel (WGS1984) \\
        WGS84\_BG & Breedtegraad van het middelpunt van het perceel (WGS1984) \\
        \bottomrule
    \end{xltabular}
}

% \begin{table}
%     \caption{Beschrijving van de attributen in de landbouwgebruikspercelen dataset \autocite{vlaanderen_landbouwpercelen_2025}}\label{app:ontwikkelingTool_landbouwkaart-vlaanderen-attributen}
%     \footnotesize
%     \begin{xltabular}{\textwidth}{@{} >{\raggedright\arraybackslash}l >{\raggedright\arraybackslash}X @{}}
%         \toprule
%         \textbf{Attribuut} & \textbf{Omschrijving} \\
%         \midrule
%         AGPAKEY & Unieke identificatiesleutel van het landbouwgebruiksperceel binnen één campagne \\
%         BT\_OMSCH & Bedrijfstypologie (economische specialisatie) \\
%         BT\_BRON & Bron van de bedrijfstypologie (jaartal berekening of aangegeven specialisatie) \\
%         PRC\_NMR & Perceelsnummer van het perceel in de verzamelaanvraag \\
%         GRAF\_OPP & Grafische oppervlakte (ha, tot 1m² nauwkeurig) \\
%         REF\_ID & Europees verplicht perceelsID. Verandert mee met de verschijningsvorm van het perceel \\
%         GWSCOD\_V & Code voorteelt \\
%         GWSNAM\_V & Naam voorteelt \\
%         GWSCOD\_H & Code hoofdteelt \\
%         GWSNAM\_H & Naam hoofdteelt \\
%         GWSGRPH\_LB & Naam groep hoofdteelt \\
%         GWSCOD\_N & Code eerste nateelt \\
%         CWSNAM\_N & Naam eerste nateelt \\
%         GWSCOD\_N2 & Code tweede nateelt \\
%         GWSNAM\_N2 & Naam tweede nateelt \\
%         AMKM & Code agromilieumaatregel \\
%         AMKM\_LB & Naam agromilieumaatregel \\
%         ECOREGELING (ECOR) & Code ecoregeling \\
%         ECOREGELING\_LB & Naam ecoregeling \\
%         BLS & Code aanplantsubsidie boslandbouwsystemen \\
%         BLS\_LB & Naam aanplantsubsidie boslandbouwsystemen \\
%         GESP\_PM & Gespecialiseerde productiemethode \\
%         GESP\_PM\_LB & Beschrijving gespecialiseerde productiemethode \\
%         BOPAKCOD & Code beheerovereenkomst VLM \\
%         BOPAKCOD\_L(B) & Naam beheerovereenkomst VLM \\
%         BIOCERT & Perceel onder bio-controle bij een bio-controleorgaan (J/N) \\
%         ERO\_NAM & Erosiecode (kleur) van het perceel \\
%         STAT\_BGV & Status Blijvend Grasland (BG) \\
%         MEERJARIG\_GRASLAND (MG) & Status Meerjarig Grasland (MG6 of hoger). Voorbeeld: “MG16 = 16de jaar grasland” \\
%         LANDBSTR & Landbouwstreek waarin het middelpunt van het perceel ligt \\
%         STAT\_AAR & Status Aardappelen, opvolgen rotatieplicht \\
%         PCT\_EKBG & Percentage van het perceel dat ecologisch kwetsbaar blijvend grasland is (range) \\
%         PCT\_WETVEE(N) & Percentage van het perceel dat wetland en/of veengebied is (range) \\
%         PRC\_GEM & Gemeente waarin het middelpunt van het perceel zich bevindt \\
%         PRC\_NIS & NIS-code van de gemeente waarin het middelpunt van het perceel zich bevindt \\
%         X\_REF & X-coördinaat van het middelpunt van het perceel (Lambert) \\
%         Y\_REF & Y-coördinaat van het middelpunt van het perceel (Lambert) \\
%         WGS84\_LG & Lengtegraad van het middelpunt van het perceel (WGS1984) \\
%         WGS84\_BG & Breedtegraad van het middelpunt van het perceel (WGS1984) \\
%         \bottomrule
%     \end{xltabular}
% \end{table}